\section{Experiments}
\label{sec:experiments}

We evaluate \textbf{L2D-SLDS} (Section~\ref{sec:generative_model}) in the
online learning-to-defer setting of Section~\ref{sec:problem_formulation}, focusing on three failure modes
of offline L2D: \emph{partial feedback}, \emph{non-stationarity}, and \emph{dynamic expert
availability}. Throughout, at each round \(t\) the router observes \((\mathbf{x}_t,\mathcal{E}_t)\),
selects \(I_t\in\mathcal{E}_t\), and then observes \(\vy_t\) and the queried prediction
\(\vhaty_{t,I_t}\) (hence the realized residual \(e_t=e_{t,I_t}\)). Unless stated otherwise, we
use squared loss and no query fees (\(\psi(e)=\lVert e\rVert_2^2\), \(\beta_k\equiv 0\)). We report
averages over multiple random seeds with standard errors.
Following standard Learning-to-Defer evaluations
\citep{mozannar2021consistent, Narasimhan, mao2023twostage, montreuil2024twostagelearningtodefermultitasklearning}, we
treat each expert as a fixed prediction rule \(f_k\) and focus on learning the \emph{router} under
partial feedback. We report results on the synthetic, Melbourne, and Jena Climate benchmarks in the main paper, and
defer the FRED experiment and additional details to Appendix~\ref{sec:experiments-details}.


\textbf{Compared methods.}
We compare our \textbf{L2D-SLDS} router under partial feedback to the following baselines.
\emph{(i) Ablation:} L2D-SLDS without the shared global factor (set \(d_g=0\)).
\emph{(ii) Contextual bandits:} LinUCB \citep{li2010contextual}, a shared-parameter LinUCB over
expert-conditioned features, NeuralUCB \citep{zhou2020neuralcontextualbanditsucbbased}, linear
Thompson sampling, and ensemble sampling \citep{lu2017ensemble}.
\emph{(iii) Non-stationary bandits:} Discounted LinUCB (D-LinUCB) \citep{russac2019weighted},
CUSUM-LinUCB \citep{liu2018change}, and GLR-LinUCB \citep{besson2022efficient}
(details in Appendix~\ref{subsection:baselines}).

\textbf{Metric.}
We report the time-averaged cumulative routing cost over horizon \(T\) (Eq.~\eqref{eq:routing_objective}). Concretely, we
compute the estimate
\(
\hat{J}(\pi) \coloneqq \frac{1}{T}\sum_{t=1}^T C_{t,I_t},
\)
where \(C_{t,I_t}\) is the realized cost of deferring to the selected expert at round \(t\). Lower is better.



\begin{figure*}[tb]
    \centering
    \includegraphics[width=0.85\textwidth]{figures/corr_new_v2}
    \caption{Regime-0 expert dependence in the synthetic transfer experiment. Each heatmap shows the
    pairwise Pearson correlation (color: \([-1,1]\)) between experts' per-round losses (experts indexed
    \(0\)--\(3\)). Top row: partial feedback (only queried losses observed).
    Columns (left-to-right) show the ground-truth correlation implied by the simulator and the
    correlations estimated by each method. L2D-SLDS best recovers the block-structured correlations
    (experts \(\{0,1\}\) vs.\ \(\{2,3\}\)). }
    \label{fig:exp_synthetic_transfer_regime_estimation}
\end{figure*}

\subsection{Synthetic: Regime-Dependent Correlation and Information Transfer}
\label{sec:exp_synthetic_transfer}

\textbf{Design goal.}
We construct a controlled routing instance in which (i) experts are \emph{correlated} in a
regime-dependent way, so that observing one expert should update beliefs about others (information
transfer; Proposition~\ref{prop:cross_update}); and (ii) one expert temporarily disappears and
re-enters.

\textbf{Environment (regimes, target, context).}
We use \(M=2\) regimes and deterministic switching in blocks of length \(L=150\) over horizon
\(T=3000\) such as $z_t \coloneqq 1 + \left\lfloor \frac{t-1}{L}\right\rfloor \bmod 2$.
The target follows a regime-dependent AR(1), and the context is the one-step lag:
\begin{equation*}
% \label{eq:exp_tri_cycle_ts}
y_t = 0.8\,y_{t-1} + d_{z_t} + \eta_t,\qquad \eta_t\sim\mathcal{N}(0,\sigma_y^2).
\end{equation*}
We set the router's context to \(x_t\coloneqq y_{t-1}\).
The regime \(z_t\) is latent to the router: the router observes only \(x_t\) (before acting) and the
single queried prediction \(\hat y_{t,I_t}\) (after acting).

\textbf{Experts and availability.}
We use \(K=4\) experts indexed \(k\in\{0,1,2,3\}\). Expert \(k=1\) is removed from the available set \(\mathcal{E}_t\)
for a contiguous interval \(t\in[2000,2500]\) and then re-enters. Each expert is a one-step forecaster
\(\hat y_{t,k}=f_k(x_t)\) with a shared slope and expert-specific intercept plus noise:
\begin{equation*}
\label{eq:exp_synth_expert_rule}
\hat y_{t,k} \coloneqq 0.8\,y_{t-1} + b_k + \varepsilon_{t,k}.
\end{equation*}
We set \((b_0,b_1,b_2,b_3)=(d_1,d_1,d_2,d_2)\), so experts \(\{0,1\}\) are well-calibrated in regime
\(z_t=1\) and experts \(\{2,3\}\) are well-calibrated in regime \(z_t=2\).

To induce \emph{regime-dependent correlation} under partial feedback, we generate the expert noises as
\[
\varepsilon_{t,k} \coloneqq s_{t,g(k)} + \tilde\varepsilon_{t,k},
\qquad
g(k)\coloneqq 1+\mathbf{1}\{k\in\{2,3\}\},
\]
with independent components \(s_{t,1},s_{t,2},(\tilde\varepsilon_{t,k})_{k}\) and regime-dependent
variances $s_{t,1}\sim\mathcal{N}(0,\sigma_{z_t,1}^2), s_{t,2}\sim\mathcal{N}(0,\sigma_{z_t,2}^2),
\tilde\varepsilon_{t,k}\sim\mathcal{N}(0,\sigma_{\mathrm{id}}^2)$, where \((\sigma_{1,1}^2,\sigma_{1,2}^2)=(\sigma_{\mathrm{hi}}^2,\sigma_{\mathrm{lo}}^2)\) and
\((\sigma_{2,1}^2,\sigma_{2,2}^2)=(\sigma_{\mathrm{lo}}^2,\sigma_{\mathrm{hi}}^2)\) with
\(\sigma_{\mathrm{hi}}^2\gg\sigma_{\mathrm{lo}}^2\). This makes experts \(\{0,1\}\) strongly
correlated in regime \(1\) and experts \(\{2,3\}\) strongly correlated in regime \(2\). We report
averaged cumulative costs in Table~\ref{tab:exp_avg_costs}.


\begin{table}[ht]
\centering
\small
\setlength{\tabcolsep}{4pt}
\caption{Averaged cumulative cost \eqref{eq:routing_objective} on synthetic experiment (Section~
\ref{sec:exp_synthetic_transfer}). Experts are evaluated only when they are available. We report mean \(\pm\) standard error across five runs. Lower is better.}
\label{tab:exp_avg_costs}
\begin{tabular}{@{}lc@{}}
\toprule
Method & Averaged Cumulative Cost \\
\midrule
\textbf{L2D-SLDS} & \(\mathbf{13.58 \pm 0.07}\)  \\
L2D-SLDS w/o \(\mathbf{g}_t\) & \(14.68 \pm 0.01\)  \\
\midrule
Shared LinUCB & \(22.96 \pm 0.01\)  \\
NeuralUCB & \(21.92 \pm 0.31\)  \\
Linear Thompson Sampling & \(22.97 \pm 0.01\) \\
Ensemble Sampling & \(23.77 \pm 0.25\) \\
LinUCB & \(22.94 \pm 0.01\)  \\
\midrule
Oracle & \(9.04\)\\
\bottomrule
\end{tabular}
\end{table}

\textbf{Model Configuration.}
We use $M = 2$ regimes with shared factor dimension $d_g = 1$ and idiosyncratic dimension $d_\alpha = 1$.
The staleness horizon for pruning is $\Delta_{\max} = 500$. We simply run a small warmup of 100 steps before
running L2D-SLDS and UCBs. For NeuralUCB we use an RNN with 16 hidden layers. For both UCBs baselines we use
$\alpha=5$ and a learning rate of $0.001$.



\textbf{Correlation recovery.}
Figure~\ref{fig:exp_synthetic_transfer_regime_estimation} compares the regime-0 loss correlation
structure. The ground truth exhibits a clear block structure: experts \(\{0,1\}\)
form one correlated group while experts \(\{2,3\}\) form another.
Under partial feedback, L2D-SLDS more consistently recovers this two-block structure than
LinUCB/NeuralUCB and the no-\(\mathbf{g}_t\) ablation. Removing \(\mathbf{g}_t\) blurs the separation
and increases cross-group correlations, consistent with reduced cross-expert information transfer.
In contrast, LinUCB/NeuralUCB yield noisier or near-uniform correlation estimates, reflecting that
purely discriminative bandit updates do not maintain a coherent joint belief over experts' latent
error processes.

\textbf{Results and Analysis.}
Table~\ref{tab:exp_avg_costs} shows that \textbf{L2D-SLDS} achieves the lowest routing cost under
partial feedback (\(13.58\pm 0.07\)), improving over LinUCB/NeuralUCB by a wide margin and also
outperforming the best fixed expert. L2D-SLDS also improves over the no-\(\mathbf{g}_t\) ablation
(\(14.68\pm 0.01\) vs.\ \(13.58\pm 0.07\)); in this synthetic correlated setting, the shared factor
provides a mechanism for updating beliefs about unqueried experts from a single queried residual
(Proposition~\ref{prop:transfer}).

\subsection{Melbourne Daily Temperatures}
\label{sec:exp_melbourne}

\textbf{Environment.}
We evaluate L2D-SLDS on the daily minimum temperature series for Melbourne \citep{daily_min_temp_melbourne_kaggle}, using the column 'minimum temperatures' as the target \(y_t\). 

\textbf{Experts and availability.}
We consider \(K=5\) fixed experts with architectures specified in the configuration:
experts 0--1 are AR models and experts 2--4 are ARIMA models. To stress-test dynamic availability, we enforce time-varying expert sets by removing expert 2 on
\(t\in[800,1200]\) and expert 3 on \(t\in[500,1500]\) (inclusive endpoints in the configuration).  


\textbf{Model Configuration.}
We use L2D-SLDS in partial-feedback mode with 4 latent regimes, shared factor dimension
\(d_g=2\), and idiosyncratic dimension \(d_\alpha=8\). The staleness horizon is \(\Delta_{\max}=2000\), the consultation
cost is \(\beta=0\). EM adaptation (offline and online) is disabled. Baselines use partial feedback: LinUCB with \(\alpha=5\) and \(L_2\) regularization \(=1.0\), and NeuralUCB with
\(\alpha=5\), hidden dimension 16, and learning rate \(10^{-3}\).


\begin{table}[ht]
\centering
\small
\setlength{\tabcolsep}{4pt}
\caption{Averaged cumulative cost \eqref{eq:routing_objective} on Melbourne.
We report the mean $\pm$ standard error over five runs; lower is better.
All routing policies are evaluated over the full horizon \(T\).}
\label{tab:exp_avg_costs_melbourne}
\begin{tabular}{@{}lc@{}}
\toprule
Method & Averaged Cumulative Cost  \\
\midrule
\textbf{L2D-SLDS} & \(\mathbf{5.70 \pm 0.01}\)  \\
L2D-SLDS w/o \(\mathbf{g}_t\) & \(5.75 \pm 0.02\) \\
\midrule
Shared LinUCB & \(5.96 \pm 0.00\)  \\
NeuralUCB & \(5.88 \pm 0.04\)  \\
Linear Thompson Sampling & \(5.93 \pm 0.02\) \\
Ensemble Sampling & \(5.93 \pm 0.02\) \\
LinUCB & \(6.10 \pm 0.03\)  \\
D-LinUCB & \(5.99 \pm 0.01\) \\
\midrule
Oracle & \(4.21\)  \\
\bottomrule
\end{tabular}
\end{table}

\textbf{Results and analysis.}
Table~\ref{tab:exp_avg_costs_melbourne} reports the averaged cumulative routing cost. Under partial feedback, L2D-SLDS attains
\(5.70 \pm 0.01\), better than all baselines including the stronger shared-parameter variants
(Shared LinUCB \(5.96\), Thompson Sampling \(5.93\), Ensemble Sampling \(5.93\)) and the non-stationary
D-LinUCB (\(5.99\)). Removing the shared factor
\(\mathbf{g}_t\) degrades performance to \(5.75 \pm 0.02\).
This gap is consistent with the role of \(\mathbf{g}_t\) under censoring: Melbourne exhibits common shocks
that affect multiple experts similarly, so a shared latent component lets a single queried residual
update beliefs about \emph{unqueried} experts via the learned cross-expert dependence.

\subsection{Jena Climate}
\label{sec:exp_jena}

\textbf{Environment.}
We evaluate on the Jena Climate dataset \citep{jena_climate}, using temperature as the scalar target
over \(T=6000\) rounds with a 13-dimensional context (meteorological covariates, target lags, and
cyclical time features). We use \(K=5\) experts (two AR, three ARIMA) with dynamic availability:
experts~2 and~3 are unavailable on disjoint intervals.

\textbf{Model configuration.}
We use \(M=4\) regimes, \(d_g=2\), \(d_\alpha=13\). EM is disabled; adaptation comes from online
filtering only. See Appendix~\ref{sec:exp_jena_appendix} for full details.

\begin{table}[ht]
\centering
\small
\setlength{\tabcolsep}{4pt}
\caption{Averaged cumulative cost on Jena Climate.
Mean \(\pm\) std.\ error over five seeds; lower is better.}
\label{tab:exp_avg_costs_jena}
\begin{tabular}{@{}lc@{}}
\toprule
Method & Averaged Cumulative Cost \\
\midrule
\textbf{L2D-SLDS} & \(\mathbf{3.61 \pm 0.00}\) \\
L2D-SLDS w/o \(\mathbf{g}_t\) & \(3.66 \pm 0.00\) \\
\midrule
Shared LinUCB & \(3.92 \pm 0.00\) \\
NeuralUCB & \(4.15 \pm 0.08\) \\
Linear Thompson Sampling & \(4.07 \pm 0.11\) \\
Ensemble Sampling & \(3.98 \pm 0.07\) \\
LinUCB & \(5.14 \pm 0.00\) \\
D-LinUCB & \(5.05 \pm 0.00\) \\
CUSUM-LinUCB & \(4.61 \pm 0.00\) \\
GLR-LinUCB & \(5.86 \pm 0.00\) \\
\midrule
Oracle & \(1.57 \pm 0.00\) \\
\bottomrule
\end{tabular}
\end{table}

\textbf{Results and analysis.}
Table~\ref{tab:exp_avg_costs_jena} confirms the ranking on a longer real-world benchmark with richer
context. L2D-SLDS (\(3.61\)) outperforms all baselines, including the strongest contextual
(Shared LinUCB, \(3.92\)) and non-stationary (CUSUM-LinUCB, \(4.61\)) variants.
The three non-stationary baselines (D-LinUCB, CUSUM-LinUCB, GLR-LinUCB) do not improve over their
stationary counterparts, suggesting that brute-force forgetting or restart mechanisms are less
effective than structured regime modeling on this benchmark. Additional runtime analysis is in
Appendix~\ref{sec:exp_jena_appendix}.

\subsection{Component Ablations}
\label{sec:ablations}

To address whether each model component is necessary, we ran three targeted ablations
(detailed results in Appendix~\ref{sec:experiments-details}).

\textbf{Regime count} (\(M\in\{1,2,4\}\), real data, five seeds).
\(M=1\) is not best on either Jena or Melbourne, confirming that regimes help. The best \(M\) is
small and dataset-dependent (\(M=4\) on Jena, \(M=2\) on Melbourne), and going from \(M=1\) to
\(M=2\) adds only one extra mode-conditioned filter per expert.

\textbf{EM estimation} (synthetic, misspecified initialization, five seeds).
Starting from deliberately perturbed parameters, offline EM with 8~iterations reduces parameter RMSE
by more than half (0.0809 \(\to\) 0.0363), confirming that the estimation machinery recovers useful
structure when the model is misspecified.

\textbf{Context-dependent transitions} (synthetic with context-dependent switching).
On a synthetic where the true transition depends on context, the learned attention-based transition
model achieves 60\% transition accuracy versus 8\% for a uniform baseline, and brings the
transition-matrix Frobenius error from 0.76 to 0.64.

